\section{Constructing the planar comparison field}\label{art:planar-appendix}

This appendix supplies the geometric ingredient used in
Theorem~\ref{art:planar-cluster}. The construction applies near a compact
regular part of one positive level. It does not assume that the BFGS iterates
converge to one point of that level.

\begin{lemma}[One field with all endpoint estimates]\label{art:planar-field}
Let $F:\R^2\to\R$ be finite and convex with minimum $F_*$, and let
$C\subset\{F=\ell\}$ be a nonempty compact set, where $\ell>F_*$.
Suppose the boundary of $\{F\le\ell\}$ contains no nontrivial segment and
every point of $C$ is non-affine-critical and has the smooth or regular
convex two-branch structure in Theorem~\ref{art:planar-cluster}.
Fix valid Armijo--weak-Wolfe constants. There exist an open neighborhood
$U\supset C$, a compact buffer $Q\supset\overline U$, a field $w$ defined
on a neighborhood of $Q$, and constants
$\eta,\kappa,\delta>0$, $C_0,C_1,L\ge0$ such that:

\begin{enumerate}
\item The field is locally Lipschitz, bounded on $Q$, and every
$g\in\partial F(p)$, $p\in Q$, satisfies $g^Tw(p)\ge\eta$.
All subgradients at a fixed point have the same pairing with $w(p)$.
\item For every eligible pair $x,z\in U$ with $\norm{s}<\delta$, its entire
segment lies in $Q$ and
\[
 |\det(s,w(x+\theta s))|\ge\kappa\norm{s}
 \quad(0\le\theta\le1),\qquad (y^Tw(z))^2/t\le C_0.
\]
\item For every such pair there are $\lambda>0$ and $\beta\in\R$ with
\[
 w(z)=\lambda w(x)+\beta s,\qquad
 |\log\lambda|\le C_1d,
 \qquad d\le1\ \Longrightarrow\ \lambda^2\le1+Ld.
\]
\end{enumerate}
All neighborhoods, constants and the field are chosen before the endpoint
pair or any BFGS matrix.
\end{lemma}
\begin{proof}
\emph{Corner motion.}
At a nonsmooth point of $C$, the local branch-value map
$x\mapsto(f_1(x),f_2(x))$ has nonsingular derivative. The inverse function
theorem gives a $C^2$ curve $p_j(t)$ characterized by
$f_1(p_j(t))=f_2(p_j(t))=t$ for $t$ near $\ell$. Its $C^1$ velocity
$v_j(t)=p_j'(t)$ satisfies
\begin{equation}\label{art:corner-velocity}
 \nabla f_1(p_j(t))^Tv_j(t)=
 \nabla f_2(p_j(t))^Tv_j(t)=1.
\end{equation}
In particular, this chart has only one nonsmooth point at level $\ell$.
A smooth point has a neighborhood with no nonsmooth points. A finite cover
of $C$ by smaller buffered charts therefore contains only finitely many
level-$\ell$ corners. After shrinking their charts, the neighborhood of
$C$, and the interval of levels, every nearby nonsmooth point lies on one
of these finitely many curves. This is a compactness consequence of the
charts, rather than an assumed finite bound on future branch changes.

\emph{Joining their directions.}
At a corner, the two active gradients generate its cone of supporting
normals. Its unit directions form a compact arc of width less than $\pi$;
by \eqref{art:corner-velocity}, every such unit normal $n$ has
$n^Tv_j(t)>0$. Normal arcs at distinct level-$\ell$ corners are disjoint:
a shared supporting normal would expose both points and their full segment
in the boundary of $\{F\le\ell\}$. The strict-contour hypothesis excludes
this. The finitely many arcs have positive angular separation.

Choose fixed smooth angular cutoffs $0\le\chi_j\le1$ with disjoint supports,
equal to one on slightly enlarged normal arcs. Shrink the level interval so
that these plateaus contain the whole moving arcs and each support still
satisfies $n^Tv_j(t)>0$ uniformly. With
$n(\theta)=(\cos\theta,\sin\theta)$ and
$T(\theta)=(-\sin\theta,\cos\theta)$, define
\begin{equation}\label{art:planar-height}
 h(t,\theta)=1+\sum_j\chi_j(\theta)
                   \bigl(n(\theta)^Tv_j(t)-1\bigr).
\end{equation}
Because at most one cutoff is nonzero, this is a convex combination of
positive quantities. Thus $h$ is bounded and bounded away from zero on the
chosen compact level interval. It is $C^1$ in $t$, smooth in $\theta$, and
its derivatives used below are bounded. No convex-support-function property
of $h$ is needed.

At a smooth point, put
$n(\theta(x))=\nabla F(x)/\norm{\nabla F(x)}$ and define
\begin{equation}\label{art:planar-field-formula}
 w(x)=h(F(x),\theta(x))n(\theta(x))
       +h_\theta(F(x),\theta(x))T(\theta(x)).
\end{equation}
All expressions are periodic, so no global choice of angle is required.
Near the $j$th switching curve the relevant cutoff equals one and all others
vanish. Consequently $h=n^Tv_j$, $h_\theta=T^Tv_j$, and
\begin{equation}\label{art:planar-corner-field}
 w(x)=v_j(F(x)).
\end{equation}
This holds on both smooth sides in a whole neighborhood of the curve.
Use the same formula on the curve itself. It provides a continuous, locally
Lipschitz extension across the gradient jump, while the field is $C^1$
elsewhere. Smaller compact buffers give uniform field and derivative bounds.
At a smooth point,
\[
 \nabla F(x)^Tw(x)=\norm{\nabla F(x)}h(F(x),\theta(x))>0;
\]
at a corner every subgradient pairs to one by
\eqref{art:corner-velocity}. On sufficiently small buffers of $C$, these
pairings have a uniform positive lower bound. One way to see the needed
uniform gradient bound away from zero is the supporting inequality
$\norm{g}\norm{x-x_*}\ge F(x)-F_*$ on the compact buffer, where $x_*$ is
any minimizer and the value gap stays positive. This proves the visibility
and common-pairing assertions.

The specific angular formula has a crucial derivative cancellation:
\begin{equation}\label{art:planar-tangent-cancellation}
 \partial_\theta(hn+h_\theta T)=(h+h_{\theta\theta})T.
\end{equation}
Thus at a smooth point the derivative of $w$ maps the tangent line of a
level curve into that same line. In the orthonormal basis $(T,n)$,
\[
 Dw=\begin{pmatrix}a&b\\0&e\end{pmatrix}
\]
with bounded coefficients. The formula $hn+h_\theta T$ is classical in
normal-angle geometry \cite{geom_Bogosel2023}; see the explicit formula in
\preprintcite[Section~2, equation~(4)]{Bog22}. Here its derivative cancellation
is used for a comparison field, without a convexity requirement on $h$.

\emph{Transversality of eligible steps.}
Choose nested neighborhoods so that $\overline U$ lies in the interior of
$Q$, with all preceding bounds valid on a larger buffer. Sufficiently close
pairs in $U$ have their segments in $Q$.
If the claimed determinant lower bound failed at arbitrarily short lengths,
compactness would give coalescing endpoints tending to $p\in Q$, normalized
steps $s/\norm{s}\to e$, and segment points at which the normalized
determinant tends to zero. Bounded subgradients and their closed graph give
endpoint gradient limits $a,b\in\partial F(p)$. Continuity of $w$ makes
$e$ parallel to $w(p)$. Positivity of the common pairing and $d>0$ force
$e=-w(p)/\norm{w(p)}$. Hence
\[
 d/\norm{s}\longrightarrow-a^Te>0,
 \qquad t/\norm{s}\longrightarrow(b-a)^Te=0,
\]
contradicting $t\ge(1-c_2)d$. This proves uniform transversality, including
every point of the segment, after reducing $\delta$.

\emph{Transport charged to decrease.}
Write $q(\theta)=x+\theta s$,
$\varphi(\theta)=F(q(\theta))$ and
$D(\theta)=\det(s,w(q(\theta)))$.
Transversality makes $D$ nonzero with constant sign on $[0,1]$.
At a smooth point the basis $(T,n)$ has determinant $-1$. Writing
$s=s_TT+s_n n$, direct calculation yields
\begin{equation}\label{art:planar-determinant-derivative}
 \det(s,Dw\,s)=s_n\bigl[(a-e)s_T+b s_n\bigr].
\end{equation}
The numerator has magnitude at most a constant times $|s_n|\norm{s}$.
Since $|D|\ge\kappa\norm{s}$ and
$\varphi'=\norm{\nabla F}\,s_n$, the quotient is bounded by a constant
times $|\varphi'|$. At a corner, \eqref{art:planar-corner-field} gives
$w\circ q=v_j\circ\varphi$. The convex function $\varphi$ has a right
derivative at every $\theta<1$, and the one-sided chain rule gives the same
bound there. Therefore, for $H(\theta)=\log|D(\theta)|$,
\[
 |H'_+(\theta)|\le A|\varphi'_+(\theta)|
 \le A\bigl(\varphi'_+(\theta)+2d\bigr),
 \qquad 0\le\theta<1.
\]
The last inequality uses monotonicity of the convex right derivative and
$\varphi'(0)=-d$.

For clarity, the one-sided comparison needed here requires no finite count
of switching intersections. If a continuous function $u$ on $[a,b]$ has
$u'_+\le0$ on $[a,b)$, then $u(b)\le u(a)$. Otherwise choose $\epsilon>0$
so that $u(b)-\epsilon b>u(a)-\epsilon a$. The continuous function
$u(t)-\epsilon t$ has a minimum at some point before $b$, where its right
derivative must be nonnegative, contradicting the strictly negative bound.
Apply this fact to both
\[
 \pm\bigl(H(\theta)-H(0)\bigr)
       -A\bigl(\varphi(\theta)-\varphi(0)+2d\theta\bigr).
\]
Both start at zero and have nonpositive right derivative. Since
$\varphi(1)\le\varphi(0)$,
\[
 |H(1)-H(0)|\le2Ad.
\]
Set $\lambda=D(1)/D(0)>0$. Because $s,w(x)$ are linearly independent,
solving in that basis gives $w(z)=\lambda w(x)+\beta s$; the preceding
bound is $|\log\lambda|\le2Ad$. For $0<d\le1$, convexity of the exponential
on $[0,1]$ gives
$\lambda^2\le e^{4Ad}\le1+(e^{4A}-1)d$.
This establishes the transport constants.

\emph{The secant correction.}
A finite buffered chart cover supplies a common Lipschitz constant for the
branch gradients. First suppose the endpoints use the same smooth branch
$f$. On a larger convex chart, with Lipschitz constant $L>0$, sufficiently
close buffered endpoints satisfy
\[
 \norm{\nabla f(z)-\nabla f(x)}^2
 \le2L\bigl[f(z)-f(x)-\nabla f(x)^T(z-x)\bigr].
\]
Indeed, subtract the tangent plane at $x$, then use the descent estimate at
$z$ for the step $-(\nabla f(z)-\nabla f(x))/L$. This step stays in the
larger chart for close endpoints, and convexity makes the resulting
function value no smaller than its value at $x$. This proves the displayed
inequality. It is the local convex-gradient estimate underlying
cocoercivity \cite{WW}; compare \preprintcite[Theorem~2.5]{WW22}.
Since the branch values and gradients
at both endpoints agree with those of $F$, the bracket is at most $t$.
Boundedness of $w$ then bounds $(y^Tw(z))^2/t$ uniformly.

For different endpoint branches $i,j$ in a two-branch chart, the segment
contains a tie point $p$. There
$(\nabla f_j(p)-\nabla f_i(p))^Tw(p)=0$.
The branch gradients and $w$ are Lipschitz, so
\[
 |y^Tw(z)|\le A_1\norm{s}.
\]
Put $a=\nabla f_i(x)$, $b=\nabla f_j(x)$ and
$r=\nabla f_j(z)-\nabla f_j(x)$, so $\norm{r}\le L\norm{s}$. Then
\[
 a^Ts=-d,\qquad b^Ts=t-d-r^Ts.
\]
The two gradient rows are uniformly independent on the smaller chart.
Their inverse bound and eligibility imply
\[
 \norm{s}\le A_2(2d+t+L\norm{s}^2)
 \le A_2\left(1+\frac{2}{1-c_2}\right)t+A_2L\norm{s}^2.
\]
For sufficiently small $\norm{s}$, absorb the last term to obtain
$t\ge a_0\norm{s}$ with $a_0>0$.
Consequently $(y^Tw(z))^2/t\le A_1^2\norm{s}/a_0$, uniformly bounded.
A finite buffered cover ensures every sufficiently close pair lies in one
of the preceding charts. Taking the common smaller distance threshold
finishes all three assertions for the same field.
\end{proof}
